% Cover letter, one page. Compiled by paper/build.py to cover_letter.pdf and cover_letter.docx.
\documentclass[11pt,a4paper]{article}
\usepackage[T1]{fontenc}
\usepackage{mathptmx}
\usepackage[margin=2.2cm]{geometry}
\usepackage{url}
\setlength{\parindent}{0pt}
\setlength{\parskip}{6pt}
\pagestyle{empty}
\begin{document}

\begin{flushright}
School of Computer and Systems Sciences\\
Jawaharlal Nehru University, New Delhi 110067, India\\
{}[CHECK: date of submission]
\end{flushright}

The Editors\\
\textit{Intelligent Oncology}

Dear Editors,

The authors submit the manuscript ``Patient-level differential privacy, membership leakage and explanation fidelity in
federated brain tumour MRI classification'' for consideration as a full-length research article.

The manuscript addresses a question that affects any multi-hospital artificial intelligence model in neuro-oncology,
namely what a formal privacy guarantee costs and what it actually protects. On a deduplicated three-source brain
tumour MRI benchmark, a mirror-aware check found 66 contaminated test images that perceptual hashing had missed.
Record-level differential privacy at $\varepsilon=8$ kept a likelihood-ratio membership attack at chance for a
balanced-accuracy cost of 0.076 in federated and 0.130 in central training, whereas a per-patient guarantee needed
$\varepsilon=64$. Without privacy the same attack separated training from test patients with an AUC of 0.805, and an
unclipped noise mechanism used in prior work still exposed patients. Grad-CAM localisation and cross-seed consistency
declined with privacy in step with accuracy.

The work fits the journal's scope in artificial intelligence for oncology imaging and builds on themes the journal has
published. Ahmed et al.\ reviewed artificial intelligence for brain tumour detection, grading and response evaluation
(Intell Oncol 2026;2(3):100071, doi:10.1016/j.intonc.2026.100071), and Jung et al.\ applied deep learning to small
brain metastases (Intell Oncol 2026;2(2):100049, doi:10.1016/j.intonc.2026.100049). Li et al.\ reviewed
cryptographic privacy-preserving medical image processing (Intell Oncol 2026;2(4):100083,
doi:10.1016/j.intonc.2026.100083), and Cai discussed transparency and privacy preservation in medical artificial
intelligence (Intell Oncol 2025;1(2):172-175, doi:10.1016/j.intonc.2025.03.006). The manuscript complements these
papers with a measured privacy, utility and explanation audit for federated neuro-oncology models.

The manuscript is original, has not been published, and is not under consideration elsewhere. All authors approved
the submission. The authors declare no competing interests, and the work received no specific funding. Code, frozen
results and figure data are available in a public repository [CHECK: repository visibility and archive DOI].

Two potential reviewers are suggested in the submission system [CHECK: names, institutional e-mail addresses, no conflicts].

Yours sincerely,

Poonguzhali N (corresponding author)\\
School of Computer and Systems Sciences, Jawaharlal Nehru University, New Delhi 110067, India\\
poonguzhali@jnu.ac.in, telephone [CHECK]

On behalf of Shashi Ranjan and Poonguzhali N

\end{document}
